% =========================================================
% CHUONG: THUC HANH TONG HOP (LABS)
% Nguon: cac notebook trong thu muc baitap/
%   Lab01.ipynb, BTTL.ipynb, Reduce.ipynb, Untitled0.ipynb, 1250080036.ipynb, BTTH2.ipynb
% =========================================================

\chapter{Thực hành tổng hợp \\ (Labs)}
\label{chap:labs}

Các chương trước trình bày từng thuật toán riêng lẻ. Chương này đặt chúng vào \textbf{quy trình làm việc thực tế} thông qua bảy bài thực hành được xây dựng từ các notebook của lớp học (thư mục \texttt{baitap/}). Mỗi bài gồm: mục tiêu, dữ liệu, các bước thực hiện kèm mã nguồn, kết quả thu được khi chạy notebook, và phần \textbf{nhận xét của giảng viên} --- nơi chỉ ra những lỗi thường gặp và những bài học phương pháp luận mà con số kết quả không tự nói ra.

\noindent Sau khi hoàn thành chương này, người học có khả năng:
\begin{itemize}
    \item Thực hiện trọn vẹn một quy trình học máy: đọc dữ liệu, tiền xử lý, khảo sát tương quan, huấn luyện nhiều mô hình, tinh chỉnh siêu tham số và đánh giá.
    \item Đọc kết quả một cách phản biện: nhận ra quá khớp, nhận ra khi dữ liệu không chứa thông tin dự đoán, nhận ra rò rỉ dữ liệu trong chuỗi thời gian.
    \item Tìm và sửa lỗi lập trình tinh vi trong cài đặt thủ công SVD/PCA.
    \item Luôn so sánh mô hình với một \textbf{mô hình cơ sở} (baseline) đơn giản.
\end{itemize}

\begin{table}[H]
\centering
\begin{tabular}{|l|p{5.6cm}|p{4.2cm}|}
\hline
\textbf{Notebook} & \textbf{Nội dung} & \textbf{Liên hệ} \\
\hline
\texttt{Lab01.ipynb} & Hồi quy giá nhà Ames Housing (Lab 1) & Chương~\ref{chap:regression}, \ref{chap:ensemble} \\ \hline
\texttt{BTTL.ipynb} & Dự đoán điểm tín dụng ngân hàng (Lab 2) & Chương~\ref{chap:validation} \\ \hline
\texttt{Reduce.ipynb} & Nội suy, RobustScaler, SVD, PCA tự cài đặt (Lab 3) & Chương~\ref{chap:preprocessing}, \ref{chap:reduce-dimension} \\ \hline
\texttt{Untitled0.ipynb} & Dự báo giá cổ phiếu VNM (Lab 4) & Chương~\ref{chap:regression} \\ \hline
\texttt{1250080036.ipynb} & Nhận dạng chữ cái viết tay (Lab 5) & Chương~\ref{chap:classify}, \ref{chap:ensemble} \\ \hline
\texttt{BTTH2.ipynb} & Chữ số viết tay (Lab 6), cảm xúc bình luận tiếng Việt (Lab 7) & Chương~\ref{chap:preprocessing}, \ref{chap:classify} \\ \hline
\texttt{KMeans}, \texttt{DBScanEx} & Ví dụ phân cụm & Đã dùng ở Chương~\ref{chap:clustering} \\ \hline
\texttt{PCAExample}, \texttt{SVDEXample} & Ví dụ giảm chiều & Đã dùng ở Chương~\ref{chap:reduce-dimension} \\ \hline
\texttt{TFIDF} & Vector hóa văn bản & Đã dùng ở Chương~\ref{chap:preprocessing} \\ \hline
\end{tabular}
\caption{Các notebook trong thư mục \texttt{baitap/} và vị trí sử dụng trong giáo trình}
\end{table}

\paragraph{Chuẩn bị môi trường.} Các bài thực hành cần thêm hai thư viện ngoài bộ cài đặt chuẩn: \texttt{mlxtend} (vẽ ma trận tương quan, ma trận biểu đồ phân tán) và \texttt{yfinance} (tải dữ liệu chứng khoán):
\begin{lstlisting}[language=bash, caption={Cài đặt thư viện cho các bài thực hành}]
pip install scikit-learn pandas matplotlib seaborn mlxtend openpyxl yfinance
\end{lstlisting}

% =========================================================
\section{Lab 1 -- Hồi quy giá nhà Ames Housing}
\label{sec:lab-ames}
% =========================================================

\paragraph{Mục tiêu.} Dự đoán giá bán \texttt{SalePrice} từ năm đặc trưng: chất lượng tổng thể (\texttt{Overall Qual}, thang 1--10), tình trạng (\texttt{Overall Cond}), diện tích sàn ở (\texttt{Gr Liv Area}), diện tích tầng hầm (\texttt{Total Bsmt SF}) và có điều hòa trung tâm hay không (\texttt{Central Air}, Y/N). So sánh hồi quy tuyến tính, hồi quy đa thức và các mô hình cây.

\subsection{Đọc dữ liệu và tiền xử lý}

\begin{lstlisting}[caption={Lab 1 -- đọc và tiền xử lý dữ liệu}]
import pandas as pd
from sklearn.preprocessing import LabelEncoder

cols = ["Overall Qual", "Overall Cond", "Gr Liv Area",
        "Central Air", "Total Bsmt SF", "SalePrice"]
df = pd.read_csv("AmesHousing.txt", usecols=cols, sep="\t")

# Central Air chi co 2 gia tri Y/N -> ma hoa 0/1
df["Central Air"] = LabelEncoder().fit_transform(df["Central Air"])

print(df.isna().sum())          # Total Bsmt SF thieu 1 gia tri
df = df.dropna()                # bo 1 dong: 2930 -> 2929 dong
\end{lstlisting}

Tham số \texttt{usecols} giúp chỉ đọc những cột cần thiết. Vì \texttt{Central Air} chỉ có hai giá trị, \texttt{LabelEncoder} cho kết quả tương đương mã hóa nhị phân và không tạo thứ tự giả (xem Chương~\ref{chap:preprocessing}). Chỉ một dòng bị thiếu trên 2930 dòng nên xóa dòng là chấp nhận được.

\subsection{Khảo sát tương quan}

\begin{lstlisting}[caption={Lab 1 -- ma trận tương quan và ma trận biểu đồ phân tán}]
import matplotlib.pyplot as plt
from mlxtend.plotting import heatmap, scatterplotmatrix

cols_con = ["Overall Qual", "Overall Cond", "Gr Liv Area", "Total Bsmt SF", "SalePrice"]
heatmap(df[cols_con].corr().values, row_names=cols_con, column_names=cols_con)
plt.show()
scatterplotmatrix(df[cols_con].values, names=cols_con, figsize=(20, 10))
plt.show()
\end{lstlisting}

\hinhchuachen{ch-labs-ames-heatmap.png}{Ma trận hệ số tương quan Pearson giữa các biến số của Ames Housing (kết quả chạy \texttt{Lab01.ipynb}).}{fig:lab-ames-heatmap}

\hinhchuachen{ch-labs-ames-scatter.png}{Ma trận biểu đồ phân tán: đường chéo là histogram của từng biến, các ô còn lại là biểu đồ phân tán từng cặp biến.}{fig:lab-ames-scatter}

Ma trận tương quan cho thấy \texttt{SalePrice} tương quan mạnh với \texttt{Overall Qual} ($r = 0{,}80$), \texttt{Gr Liv Area} ($0{,}71$) và \texttt{Total Bsmt SF} ($0{,}63$), nhưng gần như không tương quan với \texttt{Overall Cond} ($-0{,}10$). Các đặc trưng cũng tương quan với nhau (chẳng hạn \texttt{Overall Qual} và \texttt{Gr Liv Area}: $0{,}57$) --- dấu hiệu đa cộng tuyến ở mức vừa phải. Ma trận biểu đồ phân tán còn cho thấy quan hệ giữa giá và diện tích có dạng ``quạt'' (phương sai tăng theo giá trị), và một vài căn diện tích rất lớn nhưng giá thấp --- chính là các ngoại lai đã bàn ở Chương~\ref{chap:preprocessing}.

\subsection{Huấn luyện và so sánh mô hình}

\begin{lstlisting}[caption={Lab 1 -- huấn luyện các mô hình}]
from sklearn.model_selection import train_test_split, GridSearchCV
from sklearn.linear_model import LinearRegression
from sklearn.preprocessing import PolynomialFeatures
from sklearn.tree import DecisionTreeRegressor
from sklearn.ensemble import (RandomForestRegressor, AdaBoostRegressor,
                              HistGradientBoostingRegressor)

X = df.drop(columns="SalePrice")
y = df["SalePrice"]
X_train, X_test, y_train, y_test = train_test_split(X, y)   # 2196 / 733 dong

lg = LinearRegression().fit(X_train, y_train)
print(lg.score(X_train, y_train), lg.score(X_test, y_test))  # R^2 train, test

pf = PolynomialFeatures(degree=2)                 # them x_i^2 va x_i*x_j
X_train_pf = pf.fit_transform(X_train)
X_test_pf = pf.transform(X_test)                  # chi transform tap test
lg2 = LinearRegression().fit(X_train_pf, y_train)

dt = DecisionTreeRegressor().fit(X_train, y_train)
grid = GridSearchCV(DecisionTreeRegressor(), param_grid={
    "max_depth": [5, 6, 7],
    "min_samples_split": [2, 3, 4],
    "min_samples_leaf": [1, 2, 3, 4]})
grid.fit(X_train, y_train)
print(grid.best_score_, grid.best_params_)

rf = RandomForestRegressor().fit(X_train, y_train)
ada = AdaBoostRegressor().fit(X_train, y_train)
hist = HistGradientBoostingRegressor().fit(X_train, y_train)
\end{lstlisting}

\begin{table}[H]
\centering
\begin{tabular}{|l|c|c|l|}
\hline
\textbf{Mô hình} & \textbf{$R^2$ train} & \textbf{$R^2$ test} & \textbf{Nhận định} \\
\hline
Hồi quy tuyến tính & 0,758 & 0,782 & Thiếu khớp nhẹ \\ \hline
Hồi quy đa thức bậc 2 & 0,848 & 0,850 & Cải thiện rõ, không quá khớp \\ \hline
Cây quyết định (mặc định) & 0,995 & 0,769 & \textbf{Quá khớp} nặng \\ \hline
Cây quyết định (đã tinh chỉnh) & 0,877 & 0,823 & Khoảng cách train/test thu hẹp \\ \hline
Random Forest (mặc định) & 0,974 & \textbf{0,853} & Tốt nhất trên tập test \\ \hline
AdaBoost (mặc định) & 0,824 & 0,779 & Kém nhất nhóm cây \\ \hline
HistGradientBoosting (mặc định) & 0,903 & 0,839 & Tốt, nhanh \\ \hline
\end{tabular}
\caption{Kết quả Lab 1 khi chạy \texttt{Lab01.ipynb}; các con số thay đổi nhẹ giữa các lần chạy vì chia dữ liệu ngẫu nhiên}
\label{tab:lab-ames-r2}
\end{table}

Bảng~\ref{tab:lab-ames-r2} minh họa trọn vẹn đánh đổi bias -- variance (Chương~\ref{chap:validation}): mô hình tuyến tính quá đơn giản; thêm đặc trưng bậc hai giúp nắm bắt quan hệ phi tuyến (ví dụ tương tác \texttt{Overall Qual} $\times$ \texttt{Gr Liv Area}); cây quyết định không giới hạn độ sâu học thuộc tập huấn luyện ($R^2 = 0{,}995$) nhưng tổng quát kém; giới hạn độ sâu bằng \texttt{GridSearchCV} hoặc lấy trung bình nhiều cây (Random Forest) khắc phục được điều đó.

\paragraph{Sai số trên tập kiểm tra.} Notebook tính MSE của bốn mô hình cây: Random Forest $9{,}77 \times 10^8$, HistGradientBoosting $1{,}07 \times 10^9$, Decision Tree $1{,}09 \times 10^9$, AdaBoost $1{,}47 \times 10^9$. Lấy căn bậc hai để đưa về đơn vị USD: RMSE lần lượt khoảng \textbf{31\,250}, \textbf{32\,720}, \textbf{33\,050} và \textbf{38\,330} USD --- dễ diễn giải hơn nhiều so với MSE.

\subsection{Mức độ quan trọng của đặc trưng}

\hinhchuachen{ch-labs-ames-coef.png}{Hệ số hồi quy tuyến tính được vẽ như ``mức độ quan trọng'' --- một cách diễn giải \textbf{dễ gây hiểu nhầm} (xem nhận xét bên dưới).}{fig:lab-ames-coef}

\hinhchuachen{ch-labs-ames-rf-importance.png}{Mức độ quan trọng của đặc trưng theo Random Forest (\texttt{feature\_importances\_}): \texttt{Overall Qual} chiếm khoảng 65\%.}{fig:lab-ames-rf-imp}

Hồi quy tuyến tính cho hệ số lớn nhất ở \texttt{Overall Qual} ($\approx 27\,512$) và \texttt{Central Air} ($\approx 13\,066$), trong khi \texttt{Gr Liv Area} chỉ có hệ số $\approx 50$. Random Forest lại xếp hạng: \texttt{Overall Qual} (0,650), \texttt{Gr Liv Area} (0,203), \texttt{Total Bsmt SF} (0,121), \texttt{Overall Cond} (0,019), \texttt{Central Air} (0,006). Hai bảng xếp hạng mâu thuẫn ở \texttt{Central Air} --- lý do được giải thích ngay dưới đây.

\subsection{Nhận xét của giảng viên}

\begin{enumerate}
    \item \textbf{Hệ số hồi quy không phải là mức độ quan trọng.} Hệ số phụ thuộc \textbf{đơn vị đo}: \texttt{Gr Liv Area} tính bằng ft$^2$ (hàng nghìn) nên hệ số 50 USD/ft$^2$ nhân với độ lệch chuẩn khoảng 500 ft$^2$ cho ảnh hưởng $\approx 25\,000$ USD --- ngang với \texttt{Overall Qual}. \texttt{Central Air} có hệ số lớn chỉ vì biến nhị phân thay đổi ``một đơn vị'' là thay đổi toàn bộ. Muốn so sánh, phải chuẩn hóa đặc trưng bằng \texttt{StandardScaler} trước khi đọc hệ số, hoặc dùng \texttt{permutation\_importance} (Chương~\ref{chap:ensemble}).
    \item \textbf{Lỗi ghi nhãn độ đo.} Ô đánh giá cuối của \texttt{Lab01.ipynb} gán \texttt{mean\_squared\_error} cho cả cột \texttt{mean\_absolute\_error} --- hai cột vì thế giống hệt nhau. Notebook \texttt{BTTL.ipynb} đã sửa đúng lỗi này. Bài học: luôn kiểm tra lại một bảng kết quả có hai cột trùng nhau.
    \item \textbf{Dùng sai tham số tìm được.} Lưới tìm kiếm cho cây quyết định trả về \texttt{min\_\allowbreak samples\_\allowbreak split} bằng 3, nhưng mô hình huấn luyện lại dùng 4. Tệ hơn, lưới thứ hai được chạy cho \textbf{Random Forest}, nhưng tham số tìm được lại đem dùng cho \textbf{Decision Tree}. Cách an toàn là dùng trực tiếp \texttt{grid.\allowbreak best\_estimator\_} (đã được huấn luyện lại trên toàn bộ tập train), hoặc \texttt{Model(**grid.\allowbreak best\_params\_)} như \texttt{BTTL.ipynb} đã làm.
    \item \textbf{Không cố định \texttt{random\_state}.} \texttt{train\_test\_split}, \texttt{RandomForestRegressor}\ldots{} đều ngẫu nhiên; mỗi lần chạy cho kết quả khác. Hãy đặt \texttt{random\_state=42} để kết quả tái lập được, và đánh giá bằng kiểm định chéo thay vì một lần chia duy nhất.
    \item \textbf{Đánh giá mô hình trên tập test nhiều lần.} So sánh bảy mô hình trên cùng một tập test rồi chọn mô hình tốt nhất khiến tập test trở thành tập validation (Chương~\ref{chap:validation}). Nên chọn mô hình bằng kiểm định chéo trên tập train và chỉ dùng tập test một lần ở cuối.
    \item \textbf{Biến mục tiêu lệch phải.} Giá nhà lệch phải và sai số tăng theo giá; dự đoán $\log(\text{SalePrice})$ thường cho mô hình tuyến tính tốt hơn (xem nghiên cứu tình huống ở Chương~\ref{chap:preprocessing}).
\end{enumerate}

\begin{lstlisting}[caption={Lab 1 -- phiên bản đã sửa: tái lập được, dùng đúng tham số, đủ ba độ đo}]
import numpy as np
from sklearn.metrics import mean_absolute_error, mean_squared_error, r2_score
from sklearn.inspection import permutation_importance

X_train, X_test, y_train, y_test = train_test_split(X, y, test_size=0.25, random_state=42)

grid = GridSearchCV(RandomForestRegressor(random_state=42), param_grid={
    "max_depth": [7, 8, 9], "min_samples_leaf": [1, 2, 3, 4]}, cv=5)
grid.fit(X_train, y_train)
best = grid.best_estimator_                         # dung dung mo hinh vua tim

models = {"LinearRegression": LinearRegression().fit(X_train, y_train),
          "RandomForest (tuned)": best,
          "HistGB": HistGradientBoostingRegressor(random_state=42).fit(X_train, y_train)}
rows = []
for name, m in models.items():
    p = m.predict(X_test)
    rows.append({"model": name,
                 "MAE": mean_absolute_error(y_test, p),
                 "RMSE": np.sqrt(mean_squared_error(y_test, p)),
                 "R2": r2_score(y_test, p)})
print(pd.DataFrame(rows).round(3))

imp = permutation_importance(best, X_test, y_test, n_repeats=10, random_state=42)
print(pd.Series(imp.importances_mean, index=X.columns).sort_values())
\end{lstlisting}

% =========================================================
\section{Lab 2 -- Dự đoán điểm tín dụng: khi dữ liệu không có thông tin}
\label{sec:lab-bank}
% =========================================================

\paragraph{Mục tiêu.} Áp dụng đúng quy trình của Lab 1 cho tệp \texttt{Banking.xlsx}: dự đoán điểm tín dụng (\texttt{Credit score}) từ giới tính (\texttt{gender}) và số tiền đầu tư hằng tháng (\texttt{Amountinvestedmonthly}). Bài này được đưa vào giáo trình chính vì nó \textbf{thất bại} --- và thất bại theo cách rất đáng học.

\begin{lstlisting}[caption={Lab 2 -- đọc dữ liệu từ Excel và chạy cùng quy trình}]
cols = ["gender", "Amountinvestedmonthly", "Credit score"]
df = pd.read_excel("Banking.xlsx", usecols=cols)    # can thu vien openpyxl
df["gender"] = LabelEncoder().fit_transform(df["gender"])
print(df.isna().sum())        # Credit score thieu 4, Amountinvestedmonthly thieu 331
df = df.dropna()              # con 6708 dong

X = df.drop(columns="Credit score")
y = df["Credit score"]
X_train, X_test, y_train, y_test = train_test_split(X, y)   # 5031 / 1677 dong
\end{lstlisting}

\hinhchuachen{ch-labs-bank-heatmap.png}{Ma trận tương quan của Lab 2: \texttt{Credit score} gần như không tương quan với cả hai đặc trưng ($r \approx 0{,}00$ và $0{,}03$).}{fig:lab-bank-heatmap}

\begin{table}[H]
\centering
\begin{tabular}{|l|c|c|}
\hline
\textbf{Mô hình} & \textbf{$R^2$ train} & \textbf{$R^2$ test} \\
\hline
Hồi quy tuyến tính & 0,001 & 0,001 \\ \hline
Hồi quy đa thức bậc 2 & 0,001 & 0,000 \\ \hline
Cây quyết định (mặc định) & \textbf{0,957} & $\mathbf{-0{,}881}$ \\ \hline
Cây quyết định (đã tinh chỉnh) & 0,015 & $-0{,}032$ \\ \hline
Random Forest (mặc định) & 0,770 & $-0{,}297$ \\ \hline
Random Forest (đã tinh chỉnh) & 0,053 & $-0{,}007$ \\ \hline
AdaBoost & $-0{,}141$ & $-0{,}138$ \\ \hline
HistGradientBoosting & 0,062 & $-0{,}023$ \\ \hline
\end{tabular}
\caption{Kết quả Lab 2 khi chạy \texttt{BTTL.ipynb}. Sai số trên tập test: MAE $\approx 39$ điểm, MSE $\approx 2100$ (RMSE $\approx 46$ điểm) với cả bốn mô hình cây.}
\label{tab:lab-bank-r2}
\end{table}

\subsection{Nhận xét của giảng viên}

\begin{enumerate}
    \item \textbf{$R^2 \approx 0$ là kết luận, không phải lỗi.} Với hồi quy tuyến tính, $R^2$ không thể vượt xa bình phương tương quan của các đặc trưng với biến mục tiêu; ở đây $0{,}03^2 \approx 0{,}001$ --- đúng như kết quả. Hai đặc trưng này \textbf{không chứa thông tin} về điểm tín dụng. Không thuật toán nào, dù phức tạp đến đâu, ``nặn'' ra được thông tin không có trong dữ liệu --- đó chính là nguyên lý ``rác vào, rác ra'' ở Chương~\ref{chap:preprocessing}.
    \item \textbf{$R^2$ âm nghĩa là thua mô hình đoán trung bình.} Mọi $R^2$ test âm trong Bảng~\ref{tab:lab-bank-r2} cho biết mô hình còn tệ hơn việc luôn dự đoán $\bar{y}$. Hãy luôn bắt đầu bằng \texttt{DummyRegressor} (chiến lược \texttt{"mean"}) làm mô hình cơ sở; một mô hình chỉ có giá trị khi thắng được mô hình cơ sở.
    \item \textbf{Quá khớp trên nhiễu.} Cây quyết định mặc định đạt $R^2 = 0{,}957$ trên tập train nhưng $-0{,}881$ trên tập test: nó đã học thuộc \textbf{nhiễu}. Tinh chỉnh siêu tham số kéo cả hai về gần 0 --- tức là mô hình ``thừa nhận'' không học được gì, một kết quả trung thực hơn.
    \item \textbf{Dữ liệu cá nhân nhạy cảm.} Tệp \texttt{Banking.xlsx} chứa họ tên, email và số an sinh xã hội (SSN) của khách hàng. Dù chỉ đọc ba cột, việc lưu trữ và chia sẻ tệp chứa thông tin định danh cá nhân phải tuân thủ quy định bảo vệ dữ liệu (ví dụ Nghị định 13/2023/NĐ-CP của Việt Nam, GDPR của châu Âu). Trước khi đưa dữ liệu vào lớp học hay kho mã nguồn, hãy xóa hoặc ẩn danh các cột định danh.
    \item \textbf{Đạo đức khi dùng giới tính làm đặc trưng.} Dùng giới tính để chấm điểm tín dụng là phân biệt đối xử và bị cấm ở nhiều quốc gia (chẳng hạn Đạo luật Cơ hội Tín dụng Bình đẳng -- ECOA của Hoa Kỳ). Một mô hình tín dụng thực tế phải được kiểm tra \textbf{tính công bằng} (fairness) giữa các nhóm, không chỉ độ chính xác (Barocas, Hardt \& Narayanan, 2023).
    \item \textbf{Xóa dòng có thể làm lệch mẫu.} 331 dòng thiếu \texttt{Amountinvestedmonthly} bị xóa mà không kiểm tra cơ chế thiếu. Nếu những khách hàng không khai báo đầu tư có đặc điểm riêng (MNAR), mẫu còn lại sẽ bị lệch --- hãy cân nhắc các bộ điền ở mục~\ref{sec:missing}.
\end{enumerate}

\begin{lstlisting}[caption={Lab 2 -- so sánh với mô hình cơ sở}]
from sklearn.dummy import DummyRegressor
from sklearn.model_selection import cross_val_score

for name, m in {"Dummy (mean)": DummyRegressor(strategy="mean"),
                "LinearRegression": LinearRegression(),
                "RandomForest": RandomForestRegressor(max_depth=7, random_state=42)}.items():
    s = cross_val_score(m, X, y, cv=5, scoring="neg_mean_absolute_error")
    print(f"{name:18s} MAE = {-s.mean():.2f}")
# Neu MAE cua cac mo hinh khong nho hon ro rang so voi Dummy -> dac trung vo dung
\end{lstlisting}

% =========================================================
\section{Lab 3 -- Nội suy, RobustScaler và tự cài đặt SVD, PCA}
\label{sec:lab-reduce}
% =========================================================

\paragraph{Mục tiêu.} Kiểm chứng bằng tay các phép biến đổi đã học, và quan trọng hơn: học cách \textbf{phát hiện lỗi} khi tự cài đặt thuật toán bằng cách đối chiếu với thư viện.

\subsection{Nội suy tuyến tính giá trị thiếu}

Với dữ liệu có \textbf{thứ tự} (chuỗi thời gian, chuỗi đo liên tiếp), điền giá trị thiếu bằng trung bình cả cột (mục~\ref{sec:missing}) bỏ phí thông tin về vị trí. \textbf{Nội suy tuyến tính} nối thẳng hai giá trị đã biết ở hai bên khoảng trống:
\[
    x_t = x_a + \frac{t - a}{b - a}\,(x_b - x_a), \qquad a < t < b .
\]

\begin{example}[Nội suy tuyến tính]
Chuỗi $(75, 74, ?, 84, 86, ?, ?, 95)$. Ô thứ 3 nằm giữa 74 và 84: $\frac{74 + 84}{2} = 79$. Hai ô thứ 6, 7 nằm giữa 86 và 95, bước tăng $\frac{95 - 86}{3} = 3$: điền 89 và 92.
\end{example}

\begin{lstlisting}[caption={Nội suy bằng pandas}]
import numpy as np
import pandas as pd

s = pd.Series([75, 74, np.nan, 84, 86, np.nan, np.nan, 95])
print(s.interpolate().tolist())   # [75.0, 74.0, 79.0, 84.0, 86.0, 89.0, 92.0, 95.0]
\end{lstlisting}

\subsection{RobustScaler tính tay}

Với $x = (3, 7, 1, 6, 8, 4, 2, 5, 9, 11)$, sắp xếp được $(1, 2, 3, 4, 5, 6, 7, 8, 9, 11)$. NumPy tính tứ phân vị bằng nội suy tuyến tính tại vị trí $(n - 1)q$: $Q_1$ ở vị trí $9 \times 0{,}25 = 2{,}25$ nên $Q_1 = 3 + 0{,}25 \times (4 - 3) = 3{,}25$; tương tự $Q_3 = 7{,}75$; trung vị $= \frac{5 + 6}{2} = 5{,}5$. Công thức của RobustScaler cho
\[
    x' = \frac{x - \text{median}}{Q_3 - Q_1} = \frac{x - 5{,}5}{4{,}5},
\]
chẳng hạn $x = 3 \mapsto -0{,}556$ và $x = 11 \mapsto 1{,}222$ --- trùng khớp với kết quả của \texttt{RobustScaler}. Trung bình của $x$ là $5{,}6$, gần trung vị vì dữ liệu này không có ngoại lai.

\subsection{Tự cài đặt SVD: một lỗi chỉ số tinh vi}

Notebook tính SVD rút gọn cho $\mathbf{X} = \begin{pmatrix} 3 & 2 & 7 \\ 1 & 4 & 6 \end{pmatrix}$ qua trị riêng của $\mathbf{X}^\top\mathbf{X}$ (mục SVD ở Chương~\ref{chap:reduce-dimension}):
\[
    \mathbf{X}^\top\mathbf{X} = \begin{pmatrix} 10 & 10 & 27 \\ 10 & 20 & 38 \\ 27 & 38 & 85 \end{pmatrix},
    \qquad \lambda \approx (110{,}69;\ 4{,}31;\ 0),
    \qquad \sigma = \sqrt{\lambda} \approx (10{,}52;\ 2{,}08).
\]
Các giá trị suy biến khớp với \texttt{TruncatedSVD}, nhưng tọa độ chiếu thì \textbf{không}. Nguyên nhân nằm ở một dòng lệnh:

\begin{lstlisting}[caption={Lỗi sắp xếp vector riêng theo dòng thay vì theo cột}]
M = X.T @ X
eigv, eigvec = np.linalg.eigh(M)          # vector rieng nam o CAC COT
orders = np.argsort(eigv)[::-1]

sigma = eigvec[orders][:, :2]             # SAI: sap xep lai cac DONG
sigma = eigvec[:, orders][:, :2]          # DUNG: sap xep lai cac COT
reduce = X @ sigma
\end{lstlisting}

\texttt{np.linalg.eigh} trả về vector riêng theo \textbf{cột}; \texttt{eigvec[orders]} hoán vị các \textbf{dòng}, tạo ra một ma trận không còn là các vector riêng. Kết quả sai là $\begin{pmatrix} -4{,}76 & 3{,}28 \\ -5{,}95 & 0{,}82 \end{pmatrix}$; kết quả đúng là $\begin{pmatrix} -7{,}75 & 1{,}40 \\ -7{,}12 & -1{,}53 \end{pmatrix}$, trùng với \texttt{TruncatedSVD} (sai khác về dấu là bình thường, xem Chương~\ref{chap:reduce-dimension}).

\subsection{Tự cài đặt PCA: hai lỗi chồng lên nhau}

Với cùng ma trận $\mathbf{X}$, trung bình cột $\bar{\mathbf{x}} = (2;\ 3;\ 6{,}5)$ và dữ liệu đã chuẩn hóa tâm là $\pm(1;\ -1;\ 0{,}5)$. Ma trận hiệp phương sai (chia $n - 1 = 1$) có đúng một trị riêng khác 0 là $\lambda_1 = 4{,}5$, vector riêng $\mathbf{w}_1 = \left(\tfrac{2}{3};\ -\tfrac{2}{3};\ \tfrac{1}{3}\right)$, và tọa độ chiếu $z = \pm 1{,}5$ --- đúng như \texttt{PCA(n\_components=1)}.

Notebook lại cho $z = (6{,}43;\ 6{,}54)$ do hai lỗi:
\begin{enumerate}
    \item Lặp lại lỗi sắp xếp dòng \texttt{eigenvectors[orders]} như ở mục trên.
    \item Chiếu \textbf{dữ liệu gốc} \texttt{X} thay vì dữ liệu đã chuẩn hóa tâm \texttt{M}. PCA phải chiếu $\mathbf{X} - \bar{\mathbf{x}}$.
\end{enumerate}
Ngoài ra notebook chia ma trận hiệp phương sai cho $n = 2$ (được $\lambda = 2{,}25$), trong khi scikit-learn chia cho $n - 1$ (được \texttt{explained\_variance\_} $= 4{,}5$). Đây không phải lỗi, nhưng cần biết để đối chiếu đúng.

\begin{lstlisting}[caption={PCA tự cài đặt -- phiên bản đúng}]
X = np.array([[3, 2, 7], [1, 4, 6]], dtype=float)
M = X - X.mean(axis=0)                     # 1) chuan hoa tam
S = M.T @ M / (X.shape[0] - 1)             # 2) chia n - 1 nhu scikit-learn
eigenvalues, eigenvectors = np.linalg.eigh(S)
orders = np.argsort(eigenvalues)[::-1]
W = eigenvectors[:, orders[:1]]            # 3) chon theo COT
print(eigenvalues[orders][:1])             # [4.5]
print(M @ W)                               # [[ 1.5] [-1.5]]  (co the doi dau)
\end{lstlisting}

\paragraph{Bài học.} Khi tự cài đặt một thuật toán, \textbf{luôn đối chiếu với thư viện} trên một ví dụ nhỏ. Giá trị suy biến đúng không chứng minh được vector đúng; chỉ khi \textit{toàn bộ} đầu ra khớp (sai khác dấu được chấp nhận) mới có thể tin cài đặt của mình.

% =========================================================
\section{Lab 4 -- Dự báo giá cổ phiếu VNM: ảo ảnh $R^2 = 0{,}998$}
\label{sec:lab-vnm}
% =========================================================

\paragraph{Mục tiêu.} Dự báo giá đóng cửa ngày mai của cổ phiếu Vinamilk (mã \texttt{VNM}, sàn HOSE) từ giá đóng cửa của 7 phiên gần nhất, bằng hồi quy tuyến tính. Đây là cách biến bài toán chuỗi thời gian thành bài toán hồi quy bằng \textbf{đặc trưng trễ} (lag features) đã giới thiệu ở Chương~\ref{chap:regression}.

\begin{lstlisting}[caption={Lab 4 -- tải dữ liệu và tạo cửa sổ trượt 7 phiên}]
import yfinance as yf
import pandas as pd
from sklearn.linear_model import LinearRegression
from sklearn.model_selection import train_test_split

df = yf.Ticker("VNM.VN").history(period="20y")   # can ket noi Internet
close = df["Close"].values

window_size = 7
n = len(close) - window_size
data = {f"x{i}": close[i: n + i] for i in range(window_size)}   # 7 phien truoc
data["y"] = close[window_size:]                                  # phien ke tiep
df2 = pd.DataFrame(data)

X, y = df2.drop(columns="y"), df2["y"]
X_train, X_test, y_train, y_test = train_test_split(X, y)        # !!! xao tron
lg = LinearRegression().fit(X_train, y_train)
print(lg.score(X_train, y_train), lg.score(X_test, y_test))      # ~0.998  ~0.998
\end{lstlisting}

\hinhchuachen{ch-labs-vnm-close.png}{Giá đóng cửa (đã điều chỉnh) của cổ phiếu VNM theo thời gian, tải bằng \texttt{yfinance}.}{fig:lab-vnm}

Mô hình đạt $R^2 \approx 0{,}998$ trên cả tập train lẫn tập test. Một kết quả ``quá đẹp để là thật'' --- và đúng là không thật.

\subsection{Nhận xét của giảng viên}

\paragraph{1. Rò rỉ dữ liệu do xáo trộn chuỗi thời gian.} \texttt{train\_test\_split} mặc định \textbf{xáo trộn} các dòng, nên tập test gồm những ngày nằm xen giữa các ngày huấn luyện. Mô hình được ``nhìn thấy'' giá của ngày hôm trước \textit{và} ngày hôm sau một điểm test --- điều không thể có khi dự báo thật. Với chuỗi thời gian phải chia theo thứ tự: huấn luyện trên quá khứ, kiểm tra trên tương lai (\texttt{shuffle=False} hoặc \texttt{TimeSeriesSplit}).

\paragraph{2. Phải so sánh với mô hình cơ sở ngây thơ.} Mô hình ``giá ngày mai bằng giá hôm nay'' (naive forecast) không cần học gì cả. Chúng tôi chạy lại notebook với dữ liệu VNM từ 03/04/2012 đến 02/11/2026 (3773 phiên) và chia theo thời gian, 75\% đầu để huấn luyện:

\begin{table}[H]
\centering
\begin{tabular}{|l|c|}
\hline
\textbf{Cách đánh giá} & \textbf{$R^2$ trên tập test} \\
\hline
Hồi quy tuyến tính, chia \textbf{xáo trộn} (như notebook) & 0,998 \\ \hline
Hồi quy tuyến tính, chia \textbf{theo thời gian} & 0,947 \\ \hline
Mô hình ngây thơ: $\hat{y}_{t+1} = y_t$ & 0,947 \\ \hline
Hồi quy tuyến tính dự báo \textbf{tỉ suất sinh lời} ngày mai & $-0{,}001$ \\ \hline
\end{tabular}
\caption{Lab 4 đánh giá lại (dữ liệu tải ngày 02/11/2026; con số thay đổi theo thời điểm tải)}
\label{tab:lab-vnm}
\end{table}

Khi chia đúng cách, mô hình chỉ ngang bằng mô hình ngây thơ. Hệ số hồi quy giải thích vì sao: trọng số của phiên gần nhất là $0{,}988$, sáu phiên còn lại gần 0 --- mô hình đã tự ``phát hiện'' ra rằng dự báo tốt nhất là \textit{lấy giá hôm nay}.

\paragraph{3. $R^2$ cao trên giá là chỉ số gây hiểu nhầm.} Giá cổ phiếu là chuỗi có xu hướng, giá hôm nay đã gần bằng giá hôm qua, nên bất kỳ mô hình nào ``chép lại'' giá gần nhất cũng có $R^2$ rất cao. Đại lượng có ý nghĩa đầu tư là \textbf{tỉ suất sinh lời} $r_t = (y_t - y_{t-1})/y_{t-1}$. Dự báo $r_{t+1}$ từ 7 tỉ suất trước đó cho $R^2 \approx 0$: các biến động giá trong quá khứ gần như không chứa thông tin về biến động ngày mai --- phù hợp với \textbf{giả thuyết thị trường hiệu quả} dạng yếu (Fama, 1970).

\begin{lstlisting}[caption={Lab 4 -- đánh giá đúng: chia theo thời gian và so sánh với mô hình ngây thơ}]
import numpy as np
from sklearn.metrics import r2_score

cut = int(0.75 * len(df2))                       # 75% dau = qua khu
X_tr, X_te = X.iloc[:cut], X.iloc[cut:]
y_tr, y_te = y.iloc[:cut], y.iloc[cut:]

lg = LinearRegression().fit(X_tr, y_tr)
print("R2 hoi quy  :", lg.score(X_te, y_te))
print("R2 ngay tho :", r2_score(y_te, X_te["x6"]))   # x6 = gia phien gan nhat
print("He so       :", lg.coef_.round(3))

# Du bao ti suat sinh loi thay vi gia
ret = pd.Series(close).pct_change().dropna().values
m = len(ret) - window_size
Xr = np.column_stack([ret[i: m + i] for i in range(window_size)])
yr = ret[window_size:]
c = int(0.75 * len(yr))
print("R2 loi suat :", LinearRegression().fit(Xr[:c], yr[:c]).score(Xr[c:], yr[c:]))
\end{lstlisting}

\paragraph{4. Một số lưu ý khác.} Cột \texttt{Close} của \texttt{yfinance} là giá đã điều chỉnh theo cổ tức và chia tách cổ phiếu, nên giá năm 2012 thấp hơn giá giao dịch thực tế khi đó. Dữ liệu tải về thay đổi theo ngày chạy, vì vậy kết quả của mỗi người sẽ khác nhau đôi chút; hãy lưu dữ liệu ra tệp CSV để tái lập được thí nghiệm.

% =========================================================
\section{Lab 5 -- Nhận dạng chữ cái viết tay từ điểm ảnh}
\label{sec:lab-handwriting}
% =========================================================

\paragraph{Mục tiêu.} Phân loại ảnh chữ cái viết tay thường \texttt{a}--\texttt{z} bằng KNN, cây quyết định và Random Forest; đọc ảnh bằng OpenCV, biến ảnh thành vector đặc trưng, tinh chỉnh siêu tham số và đánh giá bài toán nhiều lớp.

\paragraph{Dữ liệu.} Thư mục \texttt{alpha2/} có một thư mục con cho mỗi chữ cái, tổng cộng 1144 ảnh của 26 lớp (đúng 44 ảnh mỗi lớp). Tập kiểm tra gồm 286 ảnh. Đoán ngẫu nhiên đều giữa 26 lớp chỉ đạt khoảng $1/26 \approx 3{,}8\%$; đây là mốc để so sánh các kết quả dưới đây.

\subsection{Từ ảnh đến ma trận đặc trưng}

Ảnh xám là ma trận số nguyên độ sáng. Các mô hình của chương trước nhận mỗi mẫu là một vector, nên mỗi ảnh được đưa về cùng kích thước $64 \times 64$ rồi ``duỗi phẳng'' thành vector $64 \times 64 = 4096$ chiều.

\begin{lstlisting}[caption={Lab 5 -- đọc ảnh và tạo ma trận đặc trưng}]
# pip install opencv-python
from os import listdir
import cv2
import numpy as np

X, y = [], []
for label in listdir("./alpha2"):
    for name in listdir(f"./alpha2/{label}"):
        img = cv2.imread(f"./alpha2/{label}/{name}", cv2.IMREAD_GRAYSCALE)
        img = cv2.resize(img, (64, 64))
        X.append(img.flatten())          # ma tran 64 x 64 -> vector 4096 chieu
        y.append(label)
X, y = np.array(X), np.array(y)
print(X.shape, y.shape)                  # (1144, 4096) (1144,)
\end{lstlisting}

\hinhchuachen{ch-labs-hw-sample.png}{Một ảnh chữ \texttt{a} sau khi đưa về $64 \times 64$ điểm ảnh; màu do \texttt{plt.imshow} tự gán cho độ sáng (kết quả chạy notebook).}{fig:lab-hw-sample}

\subsection{Kết quả}

\begin{table}[H]
\centering
\begin{tabular}{|l|c|c|c|c|c|}
\hline
\textbf{Mô hình} & \textbf{Train} & \textbf{Test} & \textbf{Precision} & \textbf{Recall} & \textbf{F1} \\
\hline
KNN mặc định & 0,621 & 0,315 & --- & --- & --- \\ \hline
Cây quyết định (\texttt{max\_depth=30}) & 1,000 & 0,234 & 0,228 & 0,223 & 0,218 \\ \hline
Random Forest (\texttt{max\_depth=50}) & 1,000 & 0,343 & 0,381 & 0,335 & 0,333 \\ \hline
Random Forest (500 cây, \texttt{max\_depth=20}) & 1,000 & \textbf{0,437} & 0,450 & 0,431 & 0,424 \\ \hline
\end{tabular}
\caption{Kết quả Lab 5 khi chạy notebook; precision, recall, F1 lấy trung bình macro trên tập kiểm tra. Notebook không tính ba độ đo này cho KNN thật (xem nhận xét 1)}
\label{tab:lab-hw}
\end{table}

\hinhchuachen{ch-labs-hw-cm.png}{Ma trận nhầm lẫn của ba mô hình trong notebook, từ trái sang phải: mô hình được đặt tên \texttt{knn} (thực chất là Random Forest, xem nhận xét 1), cây quyết định và Random Forest tốt nhất.}{fig:lab-hw-cm}

Mô hình tốt nhất đúng 43,7\%, gấp hơn 11 lần đoán ngẫu nhiên, nhưng vẫn còn xa mức dùng được. Ma trận nhầm lẫn bên phải cho thấy lỗi có quy luật: chữ \texttt{q} hay bị nhận thành \texttt{g} (5 ảnh), \texttt{t} thành \texttt{l} (4 ảnh), \texttt{i} thành \texttt{l} (3 ảnh), \texttt{a} thành \texttt{o} (3 ảnh) --- đều là những cặp có hình dạng gần nhau.

\subsection{Nhận xét của giảng viên}

\begin{enumerate}
    \item \textbf{Nhãn sai vì dùng lại tên biến.} Lưới tìm kiếm ở notebook chạy trên \texttt{rf}, nhưng sau đó có lệnh \texttt{knn = grid.best\_estimator\_}. Biến \texttt{knn} vì thế chứa một Random Forest, và dòng ``KNN'' trong bảng so sánh của notebook (test 0,343) thật ra không phải KNN. Dấu hiệu nhận ra là độ chính xác train bằng 1,0, khác hẳn mức 0,621 của KNN thật. Nên đặt tên biến theo đúng mô hình, và in \texttt{type(model)} trước khi lập bảng.
    \item \textbf{Chạy ô lệnh không theo thứ tự.} Ô tìm kiếm đầu tiên dùng \texttt{rf} trước ô tạo \texttt{rf}, nên notebook không chạy lại được từ đầu. Trước khi nộp bài, hãy chọn \textit{Restart and Run All}.
    \item \textbf{Hiểu sai tham số vị trí.} \texttt{RandomForestClassifier(30)} đặt \texttt{n\_estimators=30}, không phải \texttt{max\_depth=30}. Nên luôn ghi tên tham số.
    \item \textbf{Quá khớp nặng.} Mọi mô hình cây đều đạt 1,0 trên tập huấn luyện nhưng dưới 0,44 trên tập kiểm tra. Nguyên nhân chính là chỉ có 858 mẫu huấn luyện cho 4096 đặc trưng (lời nguyền số chiều, Chương~\ref{chap:reduce-dimension}), trung bình khoảng 33 mẫu mỗi lớp.
    \item \textbf{Điểm ảnh thô là đặc trưng kém.} Duỗi phẳng làm mất quan hệ không gian: cùng một nét chữ dịch sang phải vài điểm ảnh đã thành vector rất khác, nên khoảng cách Euclid của KNN không phản ánh độ giống nhau của hình dạng. Các hướng cải thiện: căn giữa và chuẩn hóa nét chữ, giảm chiều bằng PCA, dùng đặc trưng hình dạng như HOG (Dalal \& Triggs, 2005), hoặc dùng mạng tích chập (LeCun và cộng sự, 1998) kèm tăng cường dữ liệu.
    \item \textbf{Đánh giá chưa chắc chắn.} \texttt{train\_test\_split} không có \texttt{random\_state} và \texttt{stratify}, nên kết quả đổi theo mỗi lần chạy và các lớp có thể lệch. Với 286 ảnh kiểm tra, sai số chuẩn của độ chính xác quanh 0,44 là $\sqrt{0{,}44 \times 0{,}56 / 286} \approx 0{,}03$, vì vậy chênh lệch nhỏ giữa hai cấu hình chưa có ý nghĩa. Điểm kiểm định chéo (0,443) và điểm test (0,437) của mô hình tốt nhất khá gần nhau, là một dấu hiệu tốt.
\end{enumerate}

\begin{lstlisting}[caption={Lab 5 -- phiên bản đã sửa: mô hình cơ sở, PCA và kiểm định chéo phân tầng}]
from sklearn.dummy import DummyClassifier
from sklearn.decomposition import PCA
from sklearn.ensemble import RandomForestClassifier
from sklearn.model_selection import StratifiedKFold, cross_val_score
from sklearn.neighbors import KNeighborsClassifier
from sklearn.pipeline import make_pipeline
from sklearn.preprocessing import StandardScaler
from sklearn.svm import SVC

models = {
    "Dummy": DummyClassifier(strategy="most_frequent"),
    "KNN": make_pipeline(StandardScaler(), KNeighborsClassifier(n_neighbors=5)),
    "PCA + KNN": make_pipeline(StandardScaler(), PCA(n_components=40, random_state=42),
                               KNeighborsClassifier(n_neighbors=5, weights="distance")),
    "PCA + SVM": make_pipeline(StandardScaler(), PCA(n_components=40, random_state=42),
                               SVC(kernel="rbf")),
    "Random Forest": RandomForestClassifier(n_estimators=500, max_depth=20,
                                            n_jobs=-1, random_state=42),
}
cv = StratifiedKFold(n_splits=5, shuffle=True, random_state=42)
for name, model in models.items():
    s = cross_val_score(model, X, y, cv=cv, scoring="f1_macro")
    print(f"{name:14s} F1 macro = {s.mean():.3f} +/- {s.std():.3f}")
\end{lstlisting}

\noindent Chạy đoạn code này trên 1144 ảnh cho F1 macro: Dummy 0,003; KNN 0,317 ($\pm$0,023); PCA + SVM 0,330 ($\pm$0,022); PCA + KNN 0,374 ($\pm$0,015); Random Forest \textbf{0,477} ($\pm$0,041). PCA giúp KNN tốt hơn rõ rệt, nhưng Random Forest vẫn dẫn đầu; với độ lệch chuẩn cỡ 0,02, chênh lệch giữa KNN và PCA + SVM chưa có ý nghĩa.

% =========================================================
\section{Lab 6 -- Nhận dạng chữ số viết tay}
\label{sec:lab-digits}
% =========================================================

\paragraph{Mục tiêu.} Áp dụng đúng quy trình của Lab 5 cho bài toán chữ số, để thấy cùng một thuật toán cho kết quả khác hẳn khi dữ liệu thay đổi. Thư mục \texttt{data/handdigits/} gồm 42\,000 ảnh xám $28 \times 28$ của 10 chữ số, mỗi lớp có từ 3795 đến 4684 ảnh. Ảnh được duỗi phẳng thành vector 784 chiều.

\hinhchuachen{ch-labs-digits-samples.png}{Mỗi chữ số một ảnh mẫu trong thư mục \texttt{handdigits}.}{fig:lab-digits-samples}

\begin{lstlisting}[caption={Lab 6 -- đọc ảnh, chia dữ liệu phân tầng và so sánh các mô hình}]
from os import listdir
import cv2
import numpy as np
from sklearn.decomposition import PCA
from sklearn.dummy import DummyClassifier
from sklearn.ensemble import RandomForestClassifier
from sklearn.metrics import f1_score
from sklearn.model_selection import train_test_split
from sklearn.neighbors import KNeighborsClassifier
from sklearn.pipeline import make_pipeline
from sklearn.tree import DecisionTreeClassifier

X, y = [], []
for num in listdir("data/handdigits"):
    for name in listdir(f"data/handdigits/{num}"):
        img = cv2.imread(f"data/handdigits/{num}/{name}", cv2.IMREAD_GRAYSCALE)
        X.append(img.flatten())
        y.append(int(num))
X, y = np.array(X), np.array(y)                     # (42000, 784) (42000,)
X_tr, X_te, y_tr, y_te = train_test_split(X, y, test_size=0.25, stratify=y,
                                          random_state=42)
models = {
    "Dummy": DummyClassifier(strategy="most_frequent"),
    "KNN": KNeighborsClassifier(n_jobs=-1),
    "Decision Tree": DecisionTreeClassifier(max_depth=12, random_state=42),
    "Random Forest": RandomForestClassifier(n_jobs=-1, random_state=42),
    "PCA + KNN": make_pipeline(PCA(n_components=50, random_state=42),
                               KNeighborsClassifier(n_jobs=-1)),
}
for name, m in models.items():
    p = m.fit(X_tr, y_tr).predict(X_te)
    print(name, round((p == y_te).mean(), 3), round(f1_score(y_te, p, average="macro"), 3))
\end{lstlisting}

\begin{table}[H]
\centering
\begin{tabular}{|l|c|c|}
\hline
\textbf{Mô hình} & \textbf{Accuracy} & \textbf{F1 macro} \\
\hline
Dummy (luôn đoán lớp đông nhất) & 0,112 & 0,020 \\ \hline
Decision Tree (\texttt{max\_depth=12}) & 0,839 & 0,837 \\ \hline
Random Forest (mặc định) & 0,956 & 0,956 \\ \hline
KNN ($K = 5$) & 0,965 & 0,964 \\ \hline
PCA (50 thành phần) + KNN & \textbf{0,970} & \textbf{0,970} \\ \hline
\end{tabular}
\caption{Kết quả Lab 6 trên 10\,500 ảnh kiểm tra, \texttt{random\_state=42}}
\label{tab:lab-digits}
\end{table}

\hinhchuachen{ch-labs-digits-cm.png}{Ma trận nhầm lẫn của KNN trên tập kiểm tra; độ đậm theo thang log để thấy rõ các ô lỗi nhỏ.}{fig:lab-digits-cm}

KNN sai 371 trên 10\,500 ảnh. Các lỗi nhiều nhất đều hợp lý về hình dạng: 4 nhận thành 9 (24 ảnh), 8 thành 5 (22), 2 thành 7 (22), 7 thành 1 (17). Chữ số 8 có recall thấp nhất (0,918), chữ số 1 cao nhất (0,994).

\subsection{Nhận xét của giảng viên}

\begin{enumerate}
    \item \textbf{Vì sao Lab 5 chỉ đạt 0,44 mà Lab 6 đạt 0,97?} Thuật toán gần như không đổi; khác biệt nằm ở dữ liệu: Lab 6 có trung bình hơn 3000 ảnh huấn luyện mỗi lớp, so với 33 ảnh ở Lab 5, số lớp ít hơn và ảnh nhỏ hơn. Khi mô hình kém, hãy xem xét việc thu thập thêm dữ liệu trước khi đổi thuật toán.
    \item \textbf{PCA vừa nhanh hơn vừa tốt hơn.} Giảm từ 784 xuống 50 chiều loại bớt nhiễu ở các điểm ảnh ít thông tin, nên KNN cải thiện từ 0,965 lên 0,970 với chi phí tính khoảng cách nhỏ hơn nhiều.
    \item \textbf{Tìm siêu tham số xong nhưng không dùng.} Trong \texttt{BTTH2.ipynb}, \texttt{GridSearchCV} tìm được độ sâu 12 cho cây và 30 cho rừng, nhưng bảng so sánh vẫn dùng hai mô hình mặc định, nên cây quyết định có train 1,0 và test 0,840. Với cây, phải lấy \texttt{grid.best\_estimator\_}.
    \item \textbf{Chi phí của KNN.} KNN không huấn luyện gì, nhưng mỗi dự đoán phải so với 31\,500 ảnh, nên đây là bước chậm nhất của bài. Đó là đánh đổi của phương pháp học dựa trên mẫu.
\end{enumerate}

% =========================================================
\section{Lab 7 -- Phân tích cảm xúc bình luận tiếng Việt}
\label{sec:lab-sentiment}
% =========================================================

\paragraph{Mục tiêu.} Phân loại bình luận sản phẩm điện thoại thành tích cực hoặc tiêu cực, vận dụng Count và TF-IDF ở Chương~\ref{chap:preprocessing}, Naive Bayes và hồi quy logistic ở Chương~\ref{chap:classify}.

\paragraph{Dữ liệu.} Ba tệp \texttt{Train.csv}, \texttt{Dev.csv}, \texttt{Test.csv} lần lượt có 7786, 1112 và 2224 bình luận. Nhãn được suy ra từ số sao: từ 4 sao trở lên là tích cực, còn lại là tiêu cực. Tập kiểm tra có 62,8\% tích cực, nên một mô hình luôn đoán ``tích cực'' đã đạt accuracy 0,628; đó là mốc phải vượt qua.

\paragraph{Tách từ.} Tiếng Việt viết cách theo âm tiết, nên ``điện thoại'' sẽ bị đếm thành hai từ nếu không xử lý. Hàm \texttt{tokenize} của thư viện \texttt{pyvi} nối các âm tiết của một từ bằng dấu gạch dưới: \texttt{điện\_thoại}, \texttt{màn\_hình}. Từ điển thu được từ tập huấn luyện có 8681 từ.

\begin{lstlisting}[caption={Lab 7 -- tách từ, vector hóa và so sánh mô hình}]
# pip install pyvi
import pandas as pd
from pyvi.ViTokenizer import tokenize
from sklearn.feature_extraction.text import CountVectorizer, TfidfVectorizer
from sklearn.linear_model import LogisticRegression
from sklearn.metrics import f1_score, roc_auc_score
from sklearn.naive_bayes import MultinomialNB

train, dev, test = (pd.read_csv(f"data/{n}.csv") for n in ["Train", "Dev", "Test"])
for d in (train, dev, test):
    d["y"] = (d["n_star"] >= 4).astype(int)     # 1 = tich cuc, 0 = tieu cuc
    d["tokens"] = d["comment"].apply(tokenize)

for vec in [CountVectorizer(), TfidfVectorizer()]:
    A = vec.fit_transform(train["tokens"])      # fit chi tren tap huan luyen
    D, B = vec.transform(dev["tokens"]), vec.transform(test["tokens"])
    for model in [MultinomialNB(), LogisticRegression(max_iter=2000)]:
        model.fit(A, train["y"])
        p = model.predict(B)
        print(type(vec).__name__, type(model).__name__,
              round(model.score(D, dev["y"]), 3), round(model.score(B, test["y"]), 3),
              round(f1_score(test["y"], p, pos_label=0), 3),
              round(roc_auc_score(test["y"], model.predict_proba(B)[:, 1]), 3))
\end{lstlisting}

\begin{table}[H]
\centering
\begin{tabular}{|l|l|c|c|c|c|c|}
\hline
\textbf{Vector hóa} & \textbf{Mô hình} & \textbf{Train} & \textbf{Dev} & \textbf{Test} & \textbf{F1 lớp tiêu cực} & \textbf{AUC} \\
\hline
Count & Naive Bayes & 0,881 & 0,856 & 0,835 & 0,785 & 0,912 \\ \hline
Count & Random Forest & 1,000 & 0,861 & 0,845 & 0,780 & 0,911 \\ \hline
Count & Logistic Regression & 0,955 & 0,865 & 0,835 & 0,780 & 0,915 \\ \hline
TF-IDF & Naive Bayes & 0,845 & 0,802 & 0,800 & 0,668 & 0,911 \\ \hline
TF-IDF & Random Forest & 1,000 & 0,855 & 0,840 & 0,774 & 0,909 \\ \hline
TF-IDF & Logistic Regression & 0,902 & \textbf{0,871} & \textbf{0,855} & \textbf{0,805} & \textbf{0,927} \\ \hline
\end{tabular}
\caption{Kết quả Lab 7. Mô hình được chọn theo cột Dev; cột Test chỉ để báo cáo}
\label{tab:lab-sentiment}
\end{table}

\hinhchuachen{ch-labs-sentiment.png}{Trái: ma trận nhầm lẫn của mô hình tốt nhất trên tập kiểm tra. Phải: đường ROC của ba mô hình.}{fig:lab-sentiment}

Hệ số của hồi quy logistic cho biết từ nào đẩy dự đoán về mỗi phía, và kết quả rất dễ hiểu. Các từ tiêu cực mạnh nhất gồm \textit{tệ, thất\_vọng, chán, đơ, kém, chậm}; các từ tích cực mạnh nhất gồm \textit{tốt, ok, tuyệt\_vời, đẹp, mượt, trâu}. Riêng \textit{chê} lại mang trọng số tích cực vì thường đi trong cụm ``miễn chê'', cho thấy giới hạn của mô hình túi từ: nó không biết thứ tự và ngữ cảnh.

\subsection{Nhận xét của giảng viên}

\begin{enumerate}
    \item \textbf{Chấm điểm nhầm đặc trưng.} Trong \texttt{BTTH2.ipynb}, mô hình huấn luyện trên TF-IDF lại được chấm bằng \texttt{nb.score(X\_test, y\_test)}, với \texttt{X\_test} là ma trận Count. Hai ma trận cùng số cột nên không có báo lỗi, nhưng kết quả sai: notebook ghi 0,816 trong khi Naive Bayes + TF-IDF thật ra đạt 0,800. Lỗi kiểu này chỉ tránh được bằng cách gói bộ vector hóa và mô hình vào một \texttt{Pipeline}.
    \item \textbf{Accuracy chưa đủ.} Với dữ liệu lệch 63/37, Naive Bayes + TF-IDF có accuracy 0,800 nhưng chỉ tìm ra 54,1\% bình luận tiêu cực, trong khi đó là nhóm doanh nghiệp cần chú ý nhất. Vì thế bảng trên báo cáo F1 của lớp tiêu cực (Chương~\ref{chap:validation}).
    \item \textbf{Tập Dev bị bỏ quên.} Notebook không dùng \texttt{Dev.csv}, và chạy \texttt{GridSearchCV} với dữ liệu đã vector hóa trên toàn bộ tập huấn luyện. Nên dùng Dev để chọn mô hình, hoặc đặt bộ vector hóa trong pipeline khi kiểm định chéo để tránh rò rỉ.
    \item \textbf{Nhãn suy từ số sao có nhiễu.} Bình luận ``Pin kém còn lại miễn chê'' được 5 sao; 66 bình luận 5 sao trong tập huấn luyện có chứa chữ ``kém''. Bình luận 3 sao vốn trung lập cũng bị xếp thành tiêu cực. Cột \texttt{label} của dữ liệu có sẵn nhãn theo từng khía cạnh như pin, camera, là hướng phân tích chi tiết hơn.
    \item \textbf{Random Forest học thuộc.} Accuracy huấn luyện bằng 1,0 nhưng test không vượt hồi quy logistic. Với ma trận văn bản thưa và nhiều chiều, mô hình tuyến tính thường là lựa chọn tốt, rẻ và dễ diễn giải.
\end{enumerate}

% ---------------------------------------------------------
\section{Tóm tắt chương}
% ---------------------------------------------------------
\begin{itemize}
    \item Lab 1: đa thức bậc 2 và Random Forest cho $R^2 \approx 0{,}85$ trên Ames Housing; cây không giới hạn độ sâu quá khớp. Hệ số hồi quy chưa chuẩn hóa không phải là mức độ quan trọng; dùng đúng \texttt{best\_estimator\_}; cố định \texttt{random\_state}.
    \item Lab 2: khi đặc trưng không tương quan với mục tiêu, mọi mô hình đều thất bại; $R^2$ âm nghĩa là thua mô hình đoán trung bình. Luôn có mô hình cơ sở; chú ý dữ liệu cá nhân và tính công bằng.
    \item Lab 3: nội suy tuyến tính cho dữ liệu có thứ tự; khi tự cài đặt SVD/PCA, vector riêng nằm theo \textbf{cột}, PCA phải chuẩn hóa tâm, và luôn đối chiếu với thư viện.
    \item Lab 4: xáo trộn chuỗi thời gian gây rò rỉ dữ liệu; $R^2$ cao trên giá chỉ phản ánh tính liên tục của chuỗi; mô hình phải thắng được dự báo ngây thơ mới có giá trị.
    \item Lab 5: duỗi phẳng ảnh cho 4096 đặc trưng trên 858 mẫu huấn luyện nên mô hình quá khớp, tốt nhất đúng 43,7\%; luôn kiểm tra biến đang giữ mô hình nào, chạy lại notebook từ đầu và so sánh với mức đoán ngẫu nhiên.
    \item Lab 6: cùng quy trình đạt 0,97 trên chữ số nhờ có nhiều dữ liệu hơn hẳn; PCA vừa giảm chi phí vừa cải thiện KNN.
    \item Lab 7: tách từ tiếng Việt, TF-IDF + hồi quy logistic cho kết quả tốt nhất; đánh giá đúng ma trận đặc trưng, chú ý lớp tiêu cực và nhiễu của nhãn suy từ số sao.
\end{itemize}

% ---------------------------------------------------------
\section{Bài tập}
% ---------------------------------------------------------
\begin{enumerate}
    \item (Lab 1) Chạy lại Lab 1 với 5 lần chia dữ liệu khác nhau (\texttt{random\_state} = 0--4). Báo cáo trung bình và độ lệch chuẩn của $R^2$ test cho từng mô hình. Thứ hạng các mô hình có ổn định không?
    \item (Lab 1) Thay \texttt{SalePrice} bằng $\log(\text{SalePrice})$ và chuẩn hóa đặc trưng bằng \texttt{StandardScaler}. Sau đó so sánh hệ số hồi quy với mức độ quan trọng hoán vị của Random Forest, tính bằng hàm \texttt{permutation\_\allowbreak importance}.
    \item (Lab 1) Bổ sung các cột \texttt{Neighborhood} (One Hot), \texttt{Year Built}, \texttt{Garage Cars} vào pipeline ở Chương~\ref{chap:preprocessing}. $R^2$ cải thiện bao nhiêu?
    \item (Lab 2) Thay hai đặc trưng hiện tại bằng các cột khác của \texttt{Banking.xlsx} (sau khi đã loại bỏ họ tên, email, SSN). Dùng \texttt{DummyRegressor} làm mô hình cơ sở; đặc trưng nào thực sự có ích?
    \item (Lab 3) Viết hàm \texttt{my\_pca(X, k)} trả về tọa độ chiếu và tỉ lệ phương sai giải thích; kiểm tra trên tập \texttt{load\_iris} rằng kết quả trùng với \texttt{PCA(k)} (sai khác dấu từng cột).
    \item (Lab 5) Thử số thành phần PCA bằng 20, 40, 80 cho PCA + KNN, rồi liệt kê 5 cặp chữ bị nhầm nhiều nhất.
    \item (Lab 6) Vẽ độ chính xác của PCA + KNN theo số thành phần 10, 20, 50, 100, 200. Từ đâu độ chính xác không tăng nữa?
    \item (Lab 7) Gói \texttt{TfidfVectorizer} và \texttt{LogisticRegression} vào một \texttt{Pipeline}, chọn \texttt{C} và \texttt{ngram\_range} theo tập Dev. Thêm bigram có giúp mô hình hiểu ``miễn chê'' không?
    \item (Lab 4) Đánh giá mô hình VNM bằng \texttt{TimeSeriesSplit(n\_splits=5)} với các cửa sổ 3, 7, 14, 30 phiên. Lần nào thắng được mô hình ngây thơ? Lặp lại với một mã cổ phiếu khác.
\end{enumerate}

% ---------------------------------------------------------
\section*{Tài liệu tham khảo của chương}
\addcontentsline{toc}{section}{Tài liệu tham khảo của chương}
% ---------------------------------------------------------
\begin{enumerate}
    \item Barocas, S., Hardt, M., \& Narayanan, A. (2023). \textit{Fairness and Machine Learning: Limitations and Opportunities}. MIT Press.
    \item Breiman, L. (2001). Random Forests. \textit{Machine Learning}, 45(1), 5--32.
    \item De Cock, D. (2011). Ames, Iowa: Alternative to the Boston Housing Data as an End of Semester Regression Project. \textit{Journal of Statistics Education}, 19(3).
    \item Dalal, N., \& Triggs, B. (2005). Histograms of Oriented Gradients for Human Detection. \textit{Proceedings of CVPR 2005}, 886--893.
    \item Fama, E. F. (1970). Efficient Capital Markets: A Review of Theory and Empirical Work. \textit{The Journal of Finance}, 25(2), 383--417.
    \item Hyndman, R. J., \& Athanasopoulos, G. (2021). \textit{Forecasting: Principles and Practice} (3rd ed.). OTexts.
    \item LeCun, Y., Bottou, L., Bengio, Y., \& Haffner, P. (1998). Gradient-Based Learning Applied to Document Recognition. \textit{Proceedings of the IEEE}, 86(11), 2278--2324.
    \item Liu, Y. (H.) (2019). \textit{Python Machine Learning By Example} (2nd ed.), Chương 9: Stock Price Prediction with Regression Algorithms. Packt Publishing.
    \item Raschka, S., Liu, Y. (H.), \& Mirjalili, V. (2022). \textit{Machine Learning with PyTorch and Scikit-Learn}, Chương 9: Predicting Continuous Target Variables with Regression Analysis (bộ dữ liệu Ames Housing). Packt Publishing.
    \item Raschka, S. (2018). MLxtend: Providing machine learning and data science utilities and extensions to Python's scientific computing stack. \textit{Journal of Open Source Software}, 3(24), 638.
    \item Sarkar, D., Bali, R., \& Sharma, T. (2018). \textit{Practical Machine Learning with Python}, Chương 11: Forecasting Stock and Commodity Prices. Apress.
\end{enumerate}
